% ══ LaTeXEditor — HANDS-ON TUTORIAL ══════════════
%
%   Everything is in this one file. The instructions
%   are LaTeX comments — the lines starting with % —
%   so they are visible here in the editor and
%   invisible in the PDF on your right.
%
%   Work top to bottom. Nine steps, about fifteen
%   minutes. By the end you will have finished a real
%   one-page report and used every part of the editor
%   worth knowing.
%
%   Nothing here can break anything. Cmd+Z undoes any
%   change, including anything the AI writes.
%
%
% ── STEP 1 ─ Typeset it, before reading further ───
%
%   Press Cmd+Enter.
%
%   Then press it a SECOND time. The editor runs one
%   compiler pass per press, and cross-references need
%   the .aux file written by the first run. After the
%   first press the Issues panel says "Label(s) may
%   have changed" — that is expected, and the second
%   press clears it.
%
%   You should now see this document in the preview
%   pane. That is the whole loop: edit, Cmd+Enter,
%   look right.
%
% ──────────────────────────────────────────────────

\documentclass[11pt,a4paper]{article}

\usepackage[utf8]{inputenc}
\usepackage[T1]{fontenc}
\usepackage[margin=2.5cm]{geometry}
\usepackage{booktabs}
\usepackage{enumitem}

\setlength{\parindent}{0pt}
\setlength{\parskip}{0.65em}

% ── STEP 2 ─ Make it yours ────────────────────────
%
%   Replace "Your Name" below with your actual name,
%   press Cmd+Enter, and watch the title block change
%   in the preview.
%
%   That is all editing is here. Autosave writes the
%   file a second after you stop typing, so there is
%   no save step to remember — though Cmd+S works if
%   you want it.
%
% ──────────────────────────────────────────────────

\title{Q3 Regional Performance}
\author{Your Name}
\date{October 2026}

\begin{document}

\maketitle

\section{Summary}
\label{sec:summary}

% ── STEP 3 ─ Your first Cmd+K ─────────────────────
%
%   The paragraph below is badly written on purpose.
%
%   Put your cursor anywhere inside it — you do NOT
%   need to select anything — and press Cmd+K.
%
%   A bar opens underneath, and the text it is about
%   to change is tinted behind it. With nothing
%   selected it takes the whole paragraph you are
%   standing in.
%
%   Type:  make this clearer and less repetitive
%
%   Press Enter. Watch it stream, then read the diff:
%   removed lines in red, added in green, and where a
%   line was reworded rather than replaced, only the
%   WORDS that changed are tinted.
%
%   Press Cmd+Enter to accept.
%
%   Now press Cmd+Z. The edit is undone like any
%   other typing, and Cmd+Shift+Z puts it back.
%   Nothing the AI does is irreversible — which is
%   the point.
%
% ──────────────────────────────────────────────────

It is the case that the third quarter was a quarter in
which the business performed in a way that was better
than the way it performed in the second quarter, with
revenue that was up and also costs that were down, and
generally speaking the overall picture is a picture
that is positive.

% ── STEP 4 ─ Refining, instead of settling ────────
%
%   Cmd+K into the paragraph you just fixed and type:
%
%       rewrite this for a board audience
%
%   When the answer comes back, do NOT accept it.
%   Instead type:
%
%       shorter
%
%   and press Enter. That is the refine loop.
%
%   Two things worth knowing:
%
%   - Enter refines, Cmd+Enter accepts, Esc discards.
%     There is a close button on the bar too, if you
%     would rather click.
%
%   - A refinement re-runs against the ORIGINAL
%     paragraph, not against the previous answer. Ask
%     for "shorter" three times and you get three
%     attempts at the same task, not a paragraph
%     whittled away to nothing.
%
%   Accept whichever version you like best, or press
%   Esc to keep the one you already had.
%
% ──────────────────────────────────────────────────

\section{Regional Results}
\label{sec:regions}

% ── STEP 5 ─ Pasted data into a table ─────────────
%
%   Below is raw CSV, exactly as it arrives from a
%   spreadsheet. Select these four data lines — not
%   this comment — and copy them:
%
%       Region,Q2 Revenue,Q3 Revenue,Change
%       EMEA,1.24,1.81,+46%
%       APAC,0.92,1.13,+23%
%       Americas,2.10,2.04,-3%
%
%   Open the AI panel in the sidebar, choose
%   "Paste Data -> Table", paste into the box, Run.
%
%   Put your cursor on the blank line below this
%   comment, press "Insert at cursor", then Cmd+Enter.
%
%   Two things it had to get right. It escaped the %
%   signs as \% — a bare % would have commented out
%   the rest of the row. And it used booktabs rules
%   because this preamble loads booktabs; the skill is
%   given your preamble, so it can only use packages
%   you actually have.
%
% ──────────────────────────────────────────────────



% ── STEP 6 ─ Notes into prose ─────────────────────
%
%   Real reports start as bullets. Copy these four:
%
%       - EMEA growth from two enterprise renewals,
%         not new logos
%       - APAC steady, no surprises
%       - Americas down slightly, one large churn in
%         August
%       - watch: Americas pipeline thin for Q4
%
%   AI panel -> "Notes -> Prose" -> paste -> Run ->
%   Insert at cursor, on the blank line below.
%
%   Read what comes back carefully. The skill is told
%   to add no facts that were not in your notes, and
%   to leave a "% TODO:" comment where a note is too
%   vague to write from. A TODO is the skill telling
%   you it would otherwise have had to invent
%   something.
%
% ──────────────────────────────────────────────────



\section{Outlook}
\label{sec:outlook}

The Q4 forecast assumes the Americas pipeline recovers
to its Q2 level.

% ── STEP 7 ─ Breaking it, and fixing it ───────────
%
%   Time to see what happens when things go wrong.
%
%   Find the \inclduegraphics line a few lines below
%   and delete the % from the start of it, so it
%   becomes real LaTeX. Then press Cmd+Enter.
%
%   The build fails. The Issues panel fills up:
%   "Undefined control sequence" first, then a couple
%   of follow-on errors caused by it. That cascade is
%   normal — one bad command throws LaTeX off and
%   everything after it complains. Fix the FIRST error
%   and the rest usually vanish.
%
%   Click that first row: the editor jumps to the
%   offending line.
%
%   Now fix it without reading the log yourself:
%
%     1. Select the broken line.
%     2. AI panel -> "Fix Compilation Errors" -> Run.
%
%   You are never asked to paste the error. The panel
%   says "Also sends: N errors, preamble" because it
%   took them straight from the run you just did.
%   Apply, then Cmd+Enter.
%
%   (Prefer not to break it? Leave the line commented
%   and read on. But breaking things is the fastest
%   way to learn to trust a tool.)
%
% \inclduegraphics[width=0.8\textwidth]{chart.pdf}
%
% ──────────────────────────────────────────────────

% ── STEP 8 ─ Tighten, a preset worth knowing ──────
%
%   Select the whole Outlook section — the heading
%   line down to the end of its paragraph — and press
%   Cmd+K.
%
%   This time do not type anything. Click the
%   "Tighten" preset.
%
%   Tighten cuts roughly a third of the words while
%   keeping every claim, figure and citation. Reach
%   for it when a section runs half a page over.
%
%   The other two presets on that bar are "Proofread
%   Section" and "Fix Compilation Errors" — the three
%   you want most often. Everything else lives in the
%   AI panel.
%
% ──────────────────────────────────────────────────

% ── STEP 9 ─ The summary, written last ────────────
%
%   Summaries are written last and read first, so
%   this is the right moment.
%
%   Click somewhere in the document but select
%   NOTHING. Open the AI panel and choose
%   "Executive Summary", then Run.
%
%   With nothing selected the whole file is sent,
%   which is right here — a summary of one paragraph
%   would be useless. The panel tells you which it is
%   doing: it will say "whole file", not "selection".
%
%   Insert the result under the Summary heading near
%   the top, then press Cmd+Enter one last time.
%
% ──────────────────────────────────────────────────
%
%   Done. You have used:
%
%     Cmd+Enter      typeset, and the two-pass rule
%     Cmd+K          inline editing, refining,
%                    accepting, undoing
%     Issues panel   finding an error, jumping to it
%     AI panel       four skills, with your project's
%                    context attached automatically
%     Cmd+Z          undoing anything the AI wrote
%
%   Worth trying next, in your own time:
%
%   - The Outline panel jumps between sections in a
%     long file.
%   - The preview zooms with Cmd +/- while the pointer
%     is over it, or pinch on a trackpad. "Fit width"
%     re-fits when you drag the splitter.
%   - The Components panel inserts filled-in blocks —
%     booktabs tables, figures, theorems — when you
%     want a correct starting point for something
%     fiddly.
%   - The full reference is the LaTeXEditor Guide, in
%     the folder beside this one. It explains
%     everything this tutorial only showed you.
%
% ──────────────────────────────────────────────────

\end{document}
